\mychapter{Plano de Continuidade e Cronograma}{cap:plano_continuidade}\lhead{PLANO DE CONTINUIDADE}

Este capítulo organiza as atividades já concluídas, em andamento e planejadas para a continuidade da dissertação. Os dois experimentos reportados no Capítulo~\ref{cap:results} e discutidos no Capítulo~\ref{cap:conclusion} consolidam a etapa empírica inicial: estudo piloto na Paraíba e expansão experimental na Bahia. A etapa seguinte amplia o escopo para a base nacional, finaliza o treinamento do PSE, avalia sensibilidade dos parâmetros e transforma o \textit{pipeline} em um mapa operacional de recomendação de políticas.

\section{Visão Geral das Atividades Restantes}

\paragraph{Formulação do problema e revisão da literatura.}
Esta atividade consolidou a definição de política de reposição como regra de decisão parametrizada, a função objetivo baseada no custo total de inventário sob restrição de serviço, as famílias de políticas candidatas e o posicionamento em relação aos trabalhos relacionados. Ela fornece a base conceitual que sustenta o AIPE como \textit{framework} metodológico de seleção contextual.

\paragraph{Benchmarking do Experimento~1.}
O primeiro experimento funcionou como estudo piloto na Paraíba, com 15 lojas e um produto de alta intermitência. Essa etapa validou o ambiente de simulação, as métricas de avaliação e o portfólio inicial de 12 políticas, permitindo identificar comportamentos degenerados e motivar a ampliação experimental.

\paragraph{Pipeline Kedro e módulo de previsão.}
O \textit{pipeline} Kedro organiza o processamento dos dados, a extração de características, a previsão auxiliar e a rastreabilidade dos artefatos. Essa infraestrutura permite recompor resultados, comparar políticas sob as mesmas condições experimentais e integrar sinais preditivos à camada de entrada do AIPE.

\paragraph{Simulação do Experimento~2.}
O segundo experimento ampliou a avaliação para 145 séries do estado da Bahia, com cinco replicações independentes por política. Essa etapa permitiu análise por perfil operacional, validação estatística formal e comparação entre uma política única global e a seleção por perfil.

\paragraph{Validação estatística.}
A validação estatística formaliza a comparação entre políticas por meio de Wilcoxon pareado, Friedman, Nemenyi e tamanho de efeito. Essa atividade separa diferença amostral de evidência operacional mais consistente e reduz o risco de interpretar flutuações do \textit{benchmark} como dominância real.

\paragraph{Módulo \texttt{demand\_profiling} e PODs.}
O módulo \texttt{demand\_profiling} define perfis operacionais configuráveis usados para estratificação e como entrada interpretável do PSE. Os PODs não substituem a taxonomia Syntetos-Boylan; eles adicionam uma camada operacional que aproxima características estatísticas de situações de decisão em estoque.

\section{Expansão Nacional Multiestado e Multi-SKU}

A expansão nacional aplicará o \textit{pipeline} completo ao conjunto de dados nacional, com aproximadamente $10.4$ milhões de séries distribuídas em 26 estados e no Distrito Federal. O processamento por estado e por SKU ampliará a diversidade de regimes, regiões e produtos, permitindo testar se os padrões observados nos experimentos da Paraíba e da Bahia se mantêm em escala ampliada.

Essa etapa também aumenta a cobertura dos quadrantes Syntetos-Boylan e dos PODs. Com maior diversidade de séries, será possível avaliar heterogeneidade regional, diferenças entre categorias de produto e estabilidade das políticas recomendadas em contextos operacionais distintos.

\section{Treinamento e Avaliação do PSE}

Com os resultados dos experimentos executados e da expansão nacional, o PSE será treinado como componente supervisionado de recomendação. A entrada será formada pelas características operacionais das séries e pelos PODs; o rótulo será a política viável de melhor desempenho observada no \textit{benchmark}, segundo CTI sob restrição de NS.

A avaliação do PSE será feita por validação cruzada estratificada por POD e por testes fora da amostra geográfica. O objetivo é verificar se o meta-modelo aprende uma associação estável entre perfil operacional e política recomendada, sem transformar o AIPE em uma nova política de reposição ou em sistema transacional.

\section{Análise de Sensibilidade e Mapa de Recomendação}

A análise de sensibilidade testará como custos de falta e manutenção, prazo de entrega, limiares dos PODs e regimes de demanda afetam a política recomendada. Essa etapa é necessária porque uma recomendação obtida em um cenário de custo ou serviço pode deixar de ser adequada quando as condições operacionais mudam.

O resultado esperado é um mapa de recomendação por regime e perfil operacional. Esse mapa traduzirá os achados do \textit{benchmark} e do PSE em regras interpretáveis, indicando em quais condições cada família de política tende a ser mais apropriada.

\section{Otimização em Janela Móvel}

A versão atual avalia as políticas no horizonte experimental disponível para cada série. Essa escolha é adequada para comparar métodos sob um mesmo recorte temporal, mas pode gerar parametrizações sensíveis ao horizonte observado.

Como atividade futura, será testada a otimização em janelas móveis $[i,i+T]$, atualizando parâmetros à medida que novos ciclos de demanda são observados. Essa extensão reduz comportamento míope, aproxima o processo de uma operação real em \textit{rolling horizon} e permite avaliar quando uma política deve ser recalibrada.

\section{Produção Científica e Escrita da Dissertação}

Os resultados do \textit{benchmark}, da comparação entre política única e seleção por perfil, do treinamento do PSE e da expansão nacional serão candidatos a publicações adicionais. A submissão ao SBPO~2026 permanece como uma etapa de disseminação dos resultados iniciais, enquanto os artigos seguintes poderão explorar a validação nacional e o mapa de recomendação.

A escrita final consolidará resultados, limitações, validação nacional, análise de sensibilidade e discussão metodológica. A defesa deverá apresentar o AIPE como \textit{framework} de seleção contextual de políticas, com o PSE como componente supervisionado de recomendação.

\section{Cronograma Consolidado}

A Tabela~\ref{tab:cronograma} resume as atividades concluídas, em andamento e planejadas. A tabela deve ser lida em conjunto com as descrições anteriores: ela organiza a sequência temporal, mas a justificativa metodológica de cada atividade foi apresentada nas seções precedentes.

\begin{table}[h]
\centering
\small
\caption{Cronograma de atividades (referência: maio de 2026). \checkmark: concluído; $\bullet$: em andamento; $\circ$: planejado.}
\label{tab:cronograma}
\begin{tabular}{@{}p{7.2cm}cccc@{}}
\toprule
Atividade & 2024 S2 & 2025 S1 & 2025 S2 & 2026 \\
\midrule
Formulação do problema e revisão de literatura              & \checkmark & & & \\
Benchmarking Experimento~1 (15 lojas, PB, 12 políticas)     & \checkmark & \checkmark & & \\
Pipeline Kedro e módulo de previsão de demanda              & & \checkmark & \checkmark & \\
Simulação Experimento~2 (145 séries, BA, 5 replicações)     & & & \checkmark & \\
Validação estatística (Wilcoxon, Friedman, Cohen's $d$)     & & & \checkmark & \\
Módulo \texttt{demand\_profiling} e PODs                    & & & \checkmark & \\
Expansão nacional multiestado e multi-SKU                   & & & & $\bullet$ \\
Treinamento e avaliação do PSE                              & & & & $\bullet$ \\
Análise de sensibilidade e mapa de recomendação             & & & & $\circ$ \\
Otimização em janela móvel                                  & & & & $\circ$ \\
Submissão artigo SBPO~2026                                  & & & & $\circ$ \\
Submissão de artigos adicionais                             & & & & $\circ$ \\
Escrita e defesa da dissertação                             & & & & $\bullet$ \\
\bottomrule
\end{tabular}
\end{table}
